% from ega4-11-fr.tex
\oldpage[IV-3]{129}

algèbre de polynômes (donc noethérienne), et soit $\mathfrak J'_0$ l'idéal de
$B'_0$ engendré par les $T_i$ d'indice $i\leq m$~; on a alors un
$A_0$-homomorphisme surjectif $u_0:B'_0\to B_0$ transformant chaque $T_i$ en
$t_i$ ($1\leq i\leq n$), donc tel que $u_0(\mathfrak J'_0)=\mathfrak J_0$~;
cet homomorphisme permet de considérer $M_0$ comme un $B'_0$-module de type
fini. On pose alors
\[
B'_\lambda=B'_0\otimes_{A_0}A_\lambda
=A_\lambda[T_1,\ldots,T_n],\qquad
\mathfrak J'_\lambda=\mathfrak J'_0B'_\lambda,
\]
\[
B'=B'_0\otimes_{A_0}A=A[T_1,\ldots,T_n],\qquad
\mathfrak J'=\mathfrak J'_0B',
\]
de sorte que $\mathfrak J'$ est l'idéal de $B'$ engendré par les $T_i$
d'indice $i\leq m$~; en outre, il est clair que pour tout entier $k\geq0$,
on a $\mathfrak J'^kM=\mathfrak J^kM$ et
$\mathfrak J'_\lambda{}^kM_\lambda=\mathfrak J_\lambda^kM_\lambda$ pour tout
$\lambda$~; donc
$\operatorname{gr}_{\mathfrak J'}^*(M)=\operatorname{gr}_{\mathfrak J}^*(M)$
en tant que $A$-module et
$\operatorname{gr}_{\mathfrak J'_\lambda}^*(M_\lambda)
=\operatorname{gr}_{\mathfrak J_\lambda}^*(M_\lambda)$ en tant que
$A_\lambda$-module~; on peut donc substituer $B'$, $B'_\lambda$,
$\mathfrak J'$, $\mathfrak J'_\lambda$ à $B$, $B_\lambda$, $\mathfrak J$,
$\mathfrak J_\lambda$ respectivement dans la démonstration~; on notera en
outre que par construction $\operatorname{gr}_{\mathfrak J'}^*(B')$
s'identifie à $B'$ et $\operatorname{gr}_{\mathfrak J'_\lambda}^*(B'_\lambda)$
à $B'_\lambda$.

\medskip
\noindent\textbf{V)} \emph{Preuve du fait que
$\operatorname{gr}_{\mathfrak J}^*(M)$ est un
$\operatorname{gr}_{\mathfrak J}^*(B)$-module de présentation finie.} --- Nous
supposons donc désormais $A_0$ et les $A_\lambda$ noethériens, les
homomorphismes de transition $A_\lambda\to A_\mu$ injectifs,
$B_0=A_0[T_1,\ldots,T_n]$, $\mathfrak J_0$ étant engendré par les $T_i$
d'indice $i\leq m$~; l'anneau
$C_0=\operatorname{gr}_{\mathfrak J_0}^*(B_0)$ s'identifie donc à $B_0$ et
$C=\operatorname{gr}_{\mathfrak J}^*(B)$ s'identifie à $B$~; nous utiliserons
seulement tout d'abord dans ce qui suit le fait que
$C\cong C_0\otimes_{A_0}A$.

Nous aurons besoin de la variante suivante de (6.9.3)~:

\begin{lemma}[11.2.9.5]
\label{IV.11.2.9.5-fr}
Soient $A_0$ un anneau noethérien, $B_0$ une $A_0$-algèbre de type fini,
$\mathfrak J_0$ un idéal de $B_0$, $M_0$ un $B_0$-module de type fini. Il
existe alors une suite $(S_{0i})_{1\leq i\leq m}$ de sous-schémas de
$S_0=\operatorname{Spec}(A_0)$ ayant les propriétés suivantes~:
\begin{enumerate}[label=\arabic{enumi}$^\circ$]
\item Les espaces sous-jacents aux $S_{0i}$ sont deux à deux disjoints et
forment un recouvrement de $S_0$.
\item Pour chaque $i$, l'ensemble $S_{0i}$ est ouvert dans
$\bigcup_{j\geq i}S_{0j}$.
\item Chaque schéma $S_{0i}$ est affine.
\item Si $A_{0i}$ est l'anneau de $S_{0i}$ et si l'on pose
$B_{0i}=B_0\otimes_{A_0}A_{0i}$, $\mathfrak J_{0i}=\mathfrak J_0B_{0i}$,
alors
$\operatorname{gr}_{\mathfrak J_{0i}}^*(M_0\otimes_{A_0}A_{0i})$ est un
$A_{0i}$-module plat.
\end{enumerate}
\end{lemma}

On procède comme dans (6.9.3) par récurrence noethérienne, en supposant le
lemme vrai lorsqu'on y remplace $A_0$ par $A'_0=A_0/\mathfrak a_0$, où
l'idéal $\mathfrak a_0$ est tel que $V(\mathfrak a_0)\neq S_0$, $B_0$ par
$B'_0=B_0\otimes_{A_0}A'_0$, $\mathfrak J_0$ par
$\mathfrak J'_0=\mathfrak J_0B'_0$ et $M_0$ par $M_0\otimes_{A_0}A'_0$. On
considère le nilradical $\mathfrak N_0$ de $A_0$, et il suffit évidemment de
prouver le lemme en remplaçant $A_0$ par $A_0/\mathfrak N_0$ et $B_0$,
$\mathfrak J_0$, $M_0$ par les objets correspondants comme ci-dessus.
Autrement dit, on peut supposer que $A_0$ est \emph{réduit}~; d'autre part,
$\operatorname{gr}_{\mathfrak J_0}^*(B_0)$ est une $A_0$-algèbre de type fini
et $\operatorname{gr}_{\mathfrak J_0}^*(M_0)$ un
$\operatorname{gr}_{\mathfrak J_0}^*(B_0)$-module de type fini~; en vertu du
théorème de platitude générique (6.9.1), il existe un ouvert
$D(t_0)\subset S_0$ tel que si l'on pose $A_{01}=(A_0)_{t_0}$,
$\operatorname{gr}_{\mathfrak J_0}^*(M_0)\otimes_{A_0}A_{01}$ soit un
$A_{01}$-module plat~; mais puisque $A_{01}$ est un $A_0$-module plat,
$\operatorname{gr}_{\mathfrak J_0}^*(M_0)\otimes_{A_0}A_{01}$ s'identifie à
$\operatorname{gr}_{\mathfrak J_{01}}^*(M_0\otimes_{A_0}A_{01})$, où
$B_{01}=B_0\otimes_{A_0}A_{01}$ et $\mathfrak J_{01}=\mathfrak J_0B_{01}$.
Le complémentaire de $D(t_0)$ dans $S_0$ est alors de la forme
$V(\mathfrak a_0)$, et on conclut en appliquant l'hypothèse de récurrence
noethérienne.

Appliquons ce lemme à la situation actuelle, en gardant les mêmes notations~;
posons $U_{0i}=\bigcup_{j\leq i}S_{0j}$, qui est un ouvert quasi-compact de
$S_0$. Il y a donc une famille
$(f_{ik})_{1\leq i\leq m,\,1\leq k\leq n}$ d'éléments de $A_0$ tels que pour
chaque $i$, $U_{0i}$ soit réunion des $D(f_{ik})$ ($1\leq k\leq n$).

